\documentclass[11pt,twoside]{amsart}
\usepackage{amssymb, amsmath, enumerate, palatino, hyperref}
\usepackage[normalem]{ulem}
\usepackage{fullpage}
\usepackage[T1]{fontenc}
\renewcommand{\labelitemi}{\guillemotright}
\usepackage{mathrsfs}

\usepackage{tikz}
\tikzstyle{ball} = [circle,shading=ball, ball color=black,
    minimum size=1mm,inner sep=1.3pt]
\usetikzlibrary{shapes}

\theoremstyle{plain}
\newtheorem{prop}{Proposition}%[section]
\newtheorem{lemma}[prop]{Lemma}
\newtheorem{thm}[prop]{Theorem}
\newtheorem{obs}[prop]{Observation}
\newtheorem{app}[prop]{Application}
\newtheorem*{MainThm}{Main Theorem}
\newtheorem{cor}[prop]{Corollary}
\newtheorem{conj}[prop]{Conjecture}
\theoremstyle{remark}
\newtheorem{rmk}[prop]{Remark}
\newtheorem{prob}{Problem}
\newtheorem{bonus}[prop]{Bonus Problem}
\newtheorem{exc}{Exercise}
\theoremstyle{definition}
\newtheorem{ex}[prop]{Example}
\theoremstyle{definition}
\newtheorem{defn}[prop]{Definition}

\newcommand{\RR}{\mathbb{R}}
\newcommand{\ZZ}{\mathbb{Z}}
\newcommand{\CC}{\mathbb{C}}
\newcommand{\NN}{\mathbb{N}}
\newcommand{\QQ}{\mathbb{Q}}
\newcommand{\PP}{\mathbb{P}}

\newcommand{\cC}{\mathscr{C}}
\newcommand{\Ob}{\operatorname{Ob}}
\newcommand{\Mor}{\operatorname{Mor}}

\newcommand{\Set}{\operatorname{Set}}
\newcommand{\FinSet}{\operatorname{FinSet}}
\newcommand{\Vect}{\operatorname{Vect}}
\newcommand{\FinVect}{\operatorname{FinVect}}
\newcommand{\Mat}{\operatorname{Mat}}
\newcommand{\Gp}{\operatorname{Gp}}
\newcommand{\FinGp}{\operatorname{FinGp}}
\newcommand{\AbGp}{\operatorname{AbGp}}
\newcommand{\Ring}{\operatorname{Ring}}
\newcommand{\CommRing}{\operatorname{CommRing}}
\newcommand{\Field}{\operatorname{Field}}
\newcommand{\Top}{\operatorname{Top}}
\newcommand{\Aut}{\operatorname{Aut}}
\newcommand{\id}{\operatorname{id}}
\DeclareRobustCommand{\stir}{\genfrac\{\}{0pt}{}}

% For solutions.
\newenvironment{sol}{
  \medskip

  \noindent\textsc{Solution:}
}
{
  \medskip

}

\newcommand{\defeq}{:=}

\title{Math 113: Discrete Structures\\ Homework 23}
%\author{} %Remove the leading % and put your name here if using this .tex file as a template for your homework.

\begin{document}
\maketitle

\noindent
\textbf{Due:} Wednesday, April 1 at 10pm.
\bigskip

\begin{prob}
Consider the sequence $a_0,a_1,a_2,\dots$ determined by
\[a_0=0 \quad \text{and} \quad a_{n+1}=2a_n+1 \text{ for }n\geq0.\]

  \begin{enumerate}[(a)]
    \item Compute the first few members of the sequence. Make a conjecture for a closed formula for the sequence and prove it using induction.
    \item Let $A(x)=a_0+a_1x+a_2x^2+\dots$ be the generating function for this sequence. Find a closed formula for $A(x)$ using the recurrence for $(a_n)_{n=0}^\infty$.
    \item Use your answer to (b) to confirm the closed formula you found in (a).
  \end{enumerate}
\end{prob}

\begin{prob}  Recall that $\frac{1}{1-x}=1+x+x^2+x^3+\dots$, and recall the definition of multiplication of power series.
\begin{enumerate}[(a)]
 \item  What sequence corresponds to $\frac{1}{(1-x)^2}$?
 \item What sequence corresponds to $\frac{1}{(1-x)^3}$?
\end{enumerate}
\end{prob}

\begin{prob} (BONUS) Recall the recurrence for the sequence of Catalan numbers:
\[C_0=1 \quad \text{and} \quad C_{n+1}=\sum_{k=0}^n C_kC_{n-k} \text{ for } n\geq 0.\]
Find a closed formula for the generating function for the Catalan sequence.
\end{prob}

\end{document}
